\section{Related Work}
Flaky tests. A flaky test produces different outcomes across runs of the same software version \cite{luo2014empirical}. Developers spend considerable effort triaging failures that do not signal a real defect \cite{pinto2020vocabulary}, and repeated exposure erodes trust to the point that genuine failures are ignored \cite{gruber2022survey, rahman2018impact}; the cost to delivery pipelines is also documented \cite{herzig2015empirically, micco2017state}.
Detection. Techniques are dynamic, hybrid, or static, depending on whether they require running the test. Rerunning the suite is practiced at scale at Google and Microsoft \cite{lam2020study, micco2017state}, but its cost grows with the rerun budget, and even five reruns reveal only about 88\% of flaky tests \cite{lam2020understanding}. DeFlaker flags tests that fail without exercising modified code \cite{bell2018deflaker}; Shaker injects stress to expose concurrency-related flakiness \cite{silva2020shake}. For .NET, FlakeRepro \cite{leesatapornwongsa2022flakerepro} combines static and dynamic analysis to reproduce the failure of a concurrency-related flaky test: across 31 flaky tests from 22 Microsoft projects, it reproduced 26. It requires the test binaries and an observed failure, so it acts during debugging after flakiness has been found, whereas our approach predicts which tests are likely flaky from source code alone and can feed into it.
Vocabulary-based prediction. Pinto et al.\ \cite{pinto2020vocabulary} trained classifiers on the vocabulary of test code (F1 = 0.95 in Java). Subsequent work confirmed its robustness across feature sets and classifiers (FlakeFlagger \cite{alshammari2021flakeflagger}, FLAST \cite{verdecchia2021know}) and across languages (JavaScript \cite{soratto2023vocabulary}, five languages \cite{ahmad2025understanding}); Barbosa et al.\ \cite{barbosa2022test} show that root causes vary by ecosystem. Medeiros et al.\ \cite{medeiros2026vocabulary} studied Swift and found a vocabulary of asynchronous constructs and stability markers. No study covers C\#.
Positioning. This work differs from previous ones in four respects. First, the unit of analysis is the declaring method, with inheritance resolved, not the test case. Second, we compare two rerun rounds and measure how much the labels depend on them. Third, in addition to balanced cross-validation, we evaluate across repositories at the natural rate, with average precision and a ranking analysis. Fourth, we separate frequent, general markers from rare, specific ones, instead of treating the vocabulary as a single ranking.
Learned features. Pre-trained models such as CodeBERT \cite{fatima2022flakify} and prompting-based use of LLMs \cite{berndt2026can} provide mixed evidence: learned features help more in categorizing root causes than in binary prediction \cite{rahman2024quantizing}, and accuracy drops considerably under realistic imbalance \cite{rahman2025understanding}. Vocabulary features remain cheap and interpretable.